\documentclass[10pt,a4paper]{amsart}
\usepackage[margin=19mm]{geometry}
\usepackage{amsmath,amssymb,mathtools}
\usepackage[scr=rsfs,cal=euler]{mathalfa}
\usepackage{xfrac}
\usepackage[unicode,hidelinks,bookmarks=false]{hyperref}
\newcommand{\ie}{i.e.\ }
\newcommand{\cf}{cf.\ }
\newcommand{\resp}{respectively\ }
\newcommand{\Subsection}{\subsection}
\providecommand{\nobreakdash}{\nobreak\hbox{-}}
\setlength{\emergencystretch}{2em}
\multlinegap0pt
\usepackage[normalem]{ulem}
\usepackage{amssymb}
\usepackage{xcolor}
\usepackage{enumerate}
\usepackage{tikz}
\newtheorem{thm}{Theorem}
\newcommand{\C}{\mathbb{C}}
\title{Discrete averaging for discrete time dynamical systems}
\author{Vassili Gelfreich, Arturo Vieiro}
\date{Source: arXiv:2603.10737v1 (CC BY 4.0)}
\begin{document}
\maketitle

\noindent\emph{Source and licence.} Discrete averaging for discrete time dynamical systems. Version arXiv:2603.10737v1 (CC BY 4.0), distributed by arXiv under the Creative Commons Attribution 4.0 International licence, http://creativecommons.org/licenses/by/4.0/, as recorded in the arXiv licence metadata for this exact version and shown on the abstract page https://arxiv.org/abs/2603.10737v1. Full licence notice: \texttt{licenses/arxiv-2603.10737-CC-BY-4.0.txt}.\par\smallskip

\noindent\textit{Editorial scope.} Counted unit: the article's thm Thm:alaNeishtadt (source0.tex lines 659--679). The article's complete original proof of it is delivered with it (source0.tex lines 682--837, one \texttt{proof} environment of 156 lines). No mathematics is written, repaired or re-derived: the statement and the proof are the article's own lines, in the article's own notation, and no retained line is substituted or re-flowed.  Retained uncounted as the source-local context the proof needs: lines 217--220; lines 485--505; lines 633--635; lines 639--641; lines 643--648.  Nothing was omitted from any retained span, so the package carries no omission record.  No retained line contains a citation, so the wrapper carries no bibliography and no bibliography entry.  No retained line contains a figure environment, an image-inclusion command or drawing code, so no image is delivered and the package declares no asset dependency.  Editorial formatting only, in addition: an AMS \texttt{amsart} preamble that restates the article's own declarations and macros used by the retained lines, namely the environments \texttt{thm} on separate counters, together with the standard packages those lines need.  The retained environments print in this document as Theorem 1 in order of appearance, whereas the article prints the same items with the numbering of its own full document.  The complete native source of the article as distributed in the arXiv e-print for this exact version is bundled unchanged as \texttt{sources/196163/source0.tex}.
\begin{equation}\label{Eq:Xn_averaging}
X_{n}(x)=\sum_{k=-n_0}^{n-n_0}p_{nk}F^k(x)
,
\end{equation}

\begin{thm} \label{Thm:formal}
If a tangent to identity family $F_\varepsilon\in C^{n+1}\bigl(D\times [-\varepsilon_0,\varepsilon_0]\bigr)$
for some $\varepsilon_0>0$,
then there is a unique polynomial in $\varepsilon$ vector field
\(
g_\varepsilon(x)=\sum_{k=0}^{n-1}\varepsilon^k g_k(x)
\)
such that 
\[
F_\varepsilon=\Phi^1_{\varepsilon g_\varepsilon}+O(\varepsilon^{n+1})
\]
uniformly on every compact subset of $D$. Moreover, for   $1\le k\le m \le n$
\begin{equation}\label{Eq:g_k}
\frac{1}{k!}\left.\frac{\partial^k X_m}{\partial \varepsilon^{k}}\right|_{\varepsilon=0}=g_{k-1},    
\end{equation}
where $X_m$ is an interpolating vector field of order $m$ defined by \eqref{Eq:Xn_averaging},
and consequently
\begin{equation}\label{Eq:errorbound}
F_\varepsilon=\Phi^1_{X_m}+O(\varepsilon^{m+1}).
\end{equation}
\end{thm}

\begin{equation}\label{Eq:f-xi}
\|F-\operatorname{id}\|_{D_\delta}\le \epsilon.
\end{equation}

\begin{equation}\label{Eq:interpol_VF}
	X_m(x)=\sum_{k=1}^m\frac{(-1)^{k-1}}{k}\Delta_k(x),
\end{equation}

\begin{equation}\label{Eq:finite-diff}
	\Delta_0(x)=x,
	\qquad 
	\Delta_{k}(x)=\Delta_{k-1}(F(x))-\Delta_{k-1}(x)
	\quad\text{for  $k\ge 1$}.
\end{equation}

\begin{thm}\label{Thm:alaNeishtadt}
Let $\epsilon,\delta>0$ be such that $\epsilon/\delta \le c_0$ and $D_0\subset \C^d$.
Then for any map $F\in B_{\epsilon}(D_\delta)$  the interpolating vector field 
of order one satisfies $\|X_1\|_{D_\delta}\le\epsilon $ and  
\begin{equation}
   \bigl\|\Phi^1_{X_1}-F \bigr\|_{D_0} 
  \le 2\epsilon^2/\delta.
\end{equation}
The interpolating vector field of order $m\in [2,M_{\epsilon/\delta}]$ satisfies $\|X_m\|_{D_{\delta/3}}
  \le 2\epsilon $ and
\begin{equation}\label{Eq:Phi-f}
    \bigl\|\Phi^1_{X_m}-F \bigr\|_{D_0} 
  \le   
    3\epsilon  \left( \frac{6 (m-1)\epsilon}{\delta}\right)^{m} .
\end{equation}
In particular, for $\displaystyle m =M_{\epsilon/\delta} $ 
\begin{equation}\label{Eq:expbound}
\bigl \|\Phi^1_{X_{\mathrm{m}}}-F \bigr\|_{D_0}\le 
3\,\epsilon \exp\left(- c_0\, \delta /\epsilon\right).
\end{equation}
\end{thm}

\begin{proof}
We consider the family that connects $F$ and the identity map,
\[
f_\mu=(1-\mu)\operatorname{id} +\mu F
\]
where $\mu$ is a complex parameter.  
Obviously the  function $f_\mu$ is analytic in the same domain $D_\delta$ as the function $F=f_1$.

\noindent{\bf 1. The limit flow.} First we consider the
 vector field 
\[
X_{1,\mu}=f_\mu -\operatorname{id} =\mu(F-\operatorname{id} ).
\] 
Equation \eqref{Eq:f-xi} implies
\[\|X_{1,\mu}\|_{D_{\delta}}=\|\mu(F-\operatorname{id} )\|_{D_{\delta}}\le |\mu|\epsilon.
\]
For $|\mu|\le\mu_1=\delta/\epsilon$, this norm does not exceed $\delta$, the distance from
$D_0$ to the boundary of $D_\delta$. Consequently,
the time-one flow $\Phi^1_{X_{1,\mu}}(x)$ is defined 
for all $x\in D_0$ and 
\[\bigl\|\Phi^1_{X_{1,\mu}}-\operatorname{id} \bigr\|_{D_0}
\le |\mu|\epsilon.
\]
Then the triangle inequality implies that
\[
\bigl\|\Phi^1_{X_{1,\mu}}-f_\mu \bigr\|_{D_0}\le 
\bigl\|\Phi^1_{X_{1,\mu}}-\operatorname{id} \bigr\|_{D_0}+\bigl\|\operatorname{id} -f_\mu \bigr\|_{D_0}
\le 2 |\mu|\epsilon. \]
 Theorem~\ref{Thm:formal} implies that $\Phi^1_{X_{1,\mu}}-f_\mu=O(\mu^2)$. The Maximum Modulus Principle (MMP\footnote{We  use the following simple statement of Complex Analysis.
 If a function $g$ is analytic and bounded in an open disk   $|\mu|<r$ and $g^{(k)}(0)=0$ for $k=0,1,\ldots,m-1$,
 then  the maximum modulus principle implies that 
 $|g(\mu )|\le \left(|\mu|/r\right)^{m} \sup_{|\mu|< r} |g(\mu )|$.
 Of course, if the function  extends continuously onto the boundary of the disk,
 the supremum can be replaced by the maximum over $|\mu|=r$. For a vector valued function
 we apply the MMP componentwise.})
 on the disk $|\mu|\le\mu_1$ for $\mu=1$ implies that
\[
\bigl\|\Phi^1_{X_{1}}-F \bigr\|_{D_0}  
\le   \frac{ 2\epsilon\mu_1}{\mu_1^{2}}
=
 \frac{2\epsilon^2 }{\delta}.
\]
{\bf 2. Bounds for the finite differences in $D_{\delta/3}$.}
Now we consider the contribution of higher order corrections to the interpolating
vector fields.
Let $m\ge 2$ and
\[
\mu_m=\frac{2\delta}{3\epsilon (m-1)}.
\]
Suppose $|\mu|\le \mu_m$, $x_0\in D_{\delta/3}$ and $0\le k\le m-1$.
The finite induction in $k$ implies that $x_k:=f_\mu^k(x_0)\in D_{\delta}$ 
and, consequently, $|x_{k+1}-x_k|\le |\mu|\epsilon $. 
Then the definition \eqref{Eq:finite-diff}
implies that 
\[
    \Delta_{m}(x_0)=\sum_{k=0}^{m-1} (-1)^{m-k-1} \binom{m-1}{k}\Delta_1(x_k).
\]
Since $\sum_{k=0}^{m-1}\binom{m-1}{k}=2^{m-1}$ and $|\Delta_1(x_k)|\le |\mu|\epsilon$, the triangle inequality gives us
\[
	\left\| \Delta_{m}\right\|_{D_{\delta/3}}
	\le  
	2^{m-1}|\mu|\epsilon\,.
\]
We note that  $\Delta_k(x)=\Delta_{k-1}(f_\mu(x))-\Delta_{k-1}(x)$
and, since $f_\mu(x)=x+O(\mu)$, the Taylor series in $\mu$ for $\Delta_k$ start
with a higher order than the Taylor series for $\Delta_{k-1}$.
We conclude that   $\Delta_{m}(x_0)=O(\mu^{m})$.
Applying the MMP we get 
\[
\left\|\Delta_{m}\right\|_{D_{\delta/3}}
\le
2^{m-1}\mu_{m}\epsilon \,
\left(\frac{|\mu|}{\mu_{m}}\right)^{m} \,. 
\]
This bound also holds for $m=1$.

{\bf 3. Bounds for $\bigl\|X_{m,\mu}\bigr\|_{D_{\delta/3}}$ with $2\le m\le M_{\epsilon/\delta}$.}
Let the interpolating vector field $X_{m,\mu}$ be defined 
by \eqref{Eq:interpol_VF} with $f$ replaced by $f_\mu$.
We immediately conclude that 
  $X_{m,\mu}$ is analytic in $D_{\delta/3}$
for $|\mu|\le \mu_{m}$. Since $\mu_k$ are decreasing, 
we get  the following upper bound 
for  $|\mu|\le \mu_m/4$:
\[
\begin{split}
 \bigl\|X_{m,\mu}\bigr\|_{D_{\delta/3}}
&
\le\sum_{k=1}^m \frac{ \left\|\Delta_k\right\|_{D_{\delta/3}}}{k}
\le
\epsilon\sum_{k=1}^m \frac
{2^{k-1}|\mu|^k}
{k\mu_{k}^{k-1}} 
\le
\epsilon\sum_{k=1}^m \frac
{2^{k-1}|\mu|^k}
{\mu_{m}^{k-1}} 
\le
\frac{\epsilon |\mu|}{1-\frac{2|\mu|}{\mu_m}}
\le 
2\epsilon |\mu|.
\end{split}
\]
Since $\mu_m>4$ for $2\le m\le M_{\epsilon/\delta}$  we can substitute $\mu=1$ and get
\[
\|X_{m}\|_{D_{\delta/3}}\le 2 \epsilon.
\]
{\bf 4. Approximation of $F$ by the time-one map of $X_{m}$ in $D_0$.}
Recalling the definition of $\mu_m$ we observe that the inequality 
\(
2\epsilon|\mu|
 \le\frac{\delta}{3(m-1)}\le\frac{\delta}{3}
\) follows from $|\mu|\le \mu_m/4$.
Consequently, an orbit of the vector field $X_{m,\mu}$ with an initial condition  in $D_0$ remains in $D_{\delta/3}$ during one unit of time and
$$
\bigl\|\Phi^1_{X_{m,\mu}}-\operatorname{id} \bigr\|_{D_0}\le \|X_{m,\mu}\|_{D_{\delta/3}}\le 2\epsilon|\mu|\,.
$$
Recalling that  $\|f_\mu - \operatorname{id} \|_{D_0} 
 \leq |\mu| \epsilon$ and using the triangle inequality we get
$$
\bigl\|\Phi^1_{X_{m,\mu}}-f_\mu \bigr\|_{D_0}\le \bigl\|\Phi^1_{X_{m,\mu}}-\operatorname{id} \bigr\|_{D_0}+\bigl\|\operatorname{id} -f_\mu \bigr\|_{D_0} \le 
3\epsilon|\mu|.
$$
Theorem~\ref{Thm:formal} implies that
$\Phi^1_{X_{m,\mu}}$ has the same Taylor polynomial of degree $m$ in $\mu$ as the map
 $f_\mu$. Then
the MMP  based on the bound in the disk 
$|\mu|\le \mu_m/4$ can be applied with $\mu=1$ to get the desired upper bound \eqref{Eq:Phi-f}:
\[
\bigl\|\Phi^1_{X_{m}}-F \bigr\|_{D_0}  
\le   3\epsilon \left( \frac{ 4}{\mu_{m}}\right)^{m}
=
 3\epsilon\,\left(  \frac{6\epsilon (m-1)}{\delta}\right)^{m} .
\]
In order to obtain the exponentially small approximation error we note that the right-hand side 
of the last inequality depends on $m$ and takes the least value 
 near  
$m=M_{{\epsilon/\delta}}=\left\lfloor \delta/(6 \mathrm e\epsilon)\right\rfloor+1$.
Since
 \(
\frac{\delta}{6\epsilon (M_{\epsilon/\delta}-1)}
\ge \mathrm e
\)
we conclude that for this $m$
$$
\bigl\|\Phi^1_{X_m}-F \bigr\|_{D_0} 
\le 3 \epsilon \, \mathrm e^{-M_\epsilon} 
\le 3  \,
\epsilon
\exp
\left(-
\frac{\delta}{6\mathrm e \epsilon}
\right).
$$
This argument completes the proof of the theorem.
\end{proof}

\end{document}
